\documentclass[11pt,a4paper]{article}
\usepackage[margin=1in]{geometry}
\usepackage{times}
\usepackage{latexsym}
\usepackage{graphicx}
\usepackage{amsmath}
\usepackage{amssymb}
\usepackage{booktabs}
\usepackage{multirow}
\usepackage{xcolor}
\usepackage{subcaption}
\usepackage{natbib}
\usepackage{hyperref}

\title{When and Why Representation-Level Discriminability Fails as RL Reward: \\A Systematic Negative Result on Direction-Conditioned GRPO for Multi-Turn Sycophancy}

\author{Anonymous}

\begin{document}
\maketitle

\begin{abstract}
Representation Engineering has shown that linear directions in LLM activation spaces encode high-level behavioral concepts, including sycophancy. A natural hypothesis is that such directions can serve as reward signals for reinforcement learning from human feedback. We systematically test this hypothesis by extracting a sycophancy direction ($d_\text{syc}$) that achieves AUROC 0.82 in discriminating sycophantic capitulation from genuine belief revision, then deploying it as a GRPO reward signal through four continuous formulations (cosine with adaptive gate, raw cosine, recalibrated cosine, scaled cosine) and two discrete formulations (binary probe, continuous probe). All six representation-level reward variants fail: continuous cosine rewards either collapse to zero signal or actively penalize correct behavior (R1=$-0.305$ on strong-evidence samples), while the binary probe achieves SYCONScore 0.22---below both the behavioral-only baseline (R3-only: 0.33) and an undifferentiated consistency baseline (ACT: 0.301). We identify three structural causes: (1) $d_\text{syc}$ encodes \emph{direction} of belief change rather than its \emph{warrant}, conflating legitimate revision with sycophantic capitulation; (2) online RL training induces distribution shift that degrades probe accuracy from AUROC 0.82 to 0.68 by step 5000; (3) behavioral reward (NLI-based R3) is already a sufficient statistic with only 2.1\% R1-R3 disagreement. Cross-model replication on Llama-3-8B confirms these findings are structural, not model-specific. We provide actionable negative guidance for the RepE$\to$RL pipeline.
\end{abstract}

\section{Introduction}

Large language models exhibit sycophantic behavior---abandoning correct beliefs under social pressure~\cite{perez2023discovering,sharma2024towards}. This problem compounds in multi-turn dialogue, where commitment depth creates a ``double bind'': anti-sycophancy training penalizes \emph{all} belief revision, blocking both sycophantic capitulation and legitimate updating~\cite{wei2024measuring}.

Representation Engineering (RepE)~\cite{zou2023representation} offers an appealing solution: linear directions in activation space discriminate behavioral concepts, suggesting that such directions could serve as fine-grained reward signals for RL training. If we can identify a direction $d_\text{syc}$ that separates sycophantic from genuine responses in activation space, we could in principle use projections onto this direction as a reward signal that distinguishes ``the model is capitulating'' from ``the model is legitimately revising''---exactly the distinction that behavioral approaches struggle with.

We call this the \textbf{RepE$\to$RL hypothesis}: that representation-level discriminability (high AUROC in classifying sycophancy) translates into reward utility (effective signal for policy improvement via GRPO). This paper presents a systematic, pre-registered\footnote{Experiments were planned before execution; see Appendix for the full experimental protocol.} investigation of this hypothesis across six reward formulations, two model families, and multiple ablations.

\textbf{Our findings are uniformly negative.} Despite achieving AUROC 0.82 in discriminating sycophantic from genuine responses, the $d_\text{syc}$ direction fails as a reward signal in every formulation tested. The core contributions of this paper are:

\begin{enumerate}
\item \textbf{Complete evidence chain}: We document four continuous R1 failures (gate collapse, zero signal, mode collapse, sign inversion) and two discrete R1 failures (binary probe: distribution shift degradation; continuous probe: reward hacking), establishing that the failure is structural rather than a matter of tuning.

\item \textbf{Root cause analysis}: We identify the geometric mechanism---$d_\text{syc}$ encodes the \emph{direction} of belief change (toward/away from user's position) but not its \emph{warrant} (whether evidence justifies the change). Any single-direction reward inevitably conflates legitimate revision with sycophantic capitulation.

\item \textbf{Cross-model replication}: We replicate the failure pattern on Llama-3-8B (probe R1 SYCONScore 0.18 vs R3-only 0.30, $\Delta=-0.12$), confirming the finding is not model-specific.

\item \textbf{Actionable negative guidance}: Behavioral reward (NLI-based) is a sufficient statistic for multi-turn sycophancy reduction (R1-R3 disagreement $<3\%$); representation-level augmentation provides no incremental information and introduces training instabilities.
\end{enumerate}

\section{Related Work}

\paragraph{Sycophancy in LLMs.} Sycophancy---unwarranted agreement with users---has been documented across model families~\cite{perez2023discovering,sharma2024towards,wei2024measuring}. Multi-turn sycophancy compounds with commitment depth, creating accuracy degradation of up to 47\%~\cite{chen2024truth}. Training-time mitigations include consistency training~\cite{mitchell2024act}, behavioral decomposition~\cite{stengel2024pressure}, and targeted RLHF~\cite{bai2022training}.

\paragraph{Representation Engineering.} RepE~\cite{zou2023representation} established that linear directions in activation space encode high-level concepts. CCS~\cite{burns2023discovering} finds truth directions via consistency objectives. Activation Addition~\cite{turner2023activation} demonstrated inference-time steering. Latent-GRPO~\cite{zhang2026latent} uses spherical projection for reasoning correctness. CRAFT~\cite{wang2025craft} validates hidden-state GRPO for safety using prototype clustering.

\paragraph{Representation-Based Rewards.} Linear Probe Penalties~\cite{burns2024linear} use probes on reward models (not policy models). Knowledge-Prediction Gap work~\cite{kadavath2022language} documents representation-behavioral disagreement at 13--29\%. Our work is the first systematic investigation of whether representation discriminability translates to RL reward utility.

\section{Background and Method}

\subsection{Sycophancy Direction Extraction}

We extract a linear sycophancy direction $d_\text{syc}$ from contrastive pairs of multi-turn dialogues. Each pair consists of a \emph{sycophantic} response (model capitulates to invalid pressure) and a \emph{genuine} response (model maintains or legitimately revises belief). Pairs are stratified by evidence strength (weak/medium/strong) and style-controlled to isolate cognitive signal from surface features.

Given hidden states $h_\text{syc}^{(l)}$ and $h_\text{gen}^{(l)}$ at layer $l$ for each pair, we compute:
\begin{equation}
d_\text{syc}^{(l)} = \frac{1}{N}\sum_{i=1}^{N}(h_{\text{syc},i}^{(l)} - h_{\text{gen},i}^{(l)})
\end{equation}

The direction is extracted at the best-discriminating layer (Layer 19 for Qwen3-8B, Layer 17 for Llama-3-8B), selected by cross-validated AUROC.

\subsection{DC-GRPO Framework}

Direction-Conditioned GRPO (DC-GRPO) augments standard GRPO with a representation-level reward component:
\begin{equation}
R = w_1 \cdot R_1 + w_3 \cdot R_3
\end{equation}
where $R_1$ is the representation-level reward (multiple formulations tested) and $R_3$ is a behavioral reward based on NLI consistency checking.

\paragraph{R1 Formulations Tested:}
\begin{itemize}
\item \textbf{Cosine + tanh gate}: $R_1 = g(t) \cdot \cos(h, d_\text{syc})$ with learnable gate $g(t)$
\item \textbf{Raw cosine}: $R_1 = \cos(h, d_\text{syc})$
\item \textbf{Recalibrated cosine}: Raw cosine with $d_\text{syc}$ recomputed every 500 steps
\item \textbf{Scaled cosine}: $R_1 = 5 \cdot \cos(h, d_\text{syc})$
\item \textbf{Binary probe}: $R_1 = \text{probe}(h; d_\text{syc}) \in \{0, 1\}$
\item \textbf{Continuous probe}: $R_1 = P(\text{genuine} \mid h; d_\text{syc}) \in [0, 1]$
\end{itemize}

\paragraph{R3 Behavioral Reward.} $R_3$ uses NLI to assess whether the model's response is consistent with available evidence. It directly evaluates output quality without requiring access to internal representations.

\subsection{Evaluation Metrics}

\paragraph{SYCONScore.} Harmonic mean of Truth Flip rate (ToF: correctly revising beliefs given valid evidence) and No-Flip rate (NoF: maintaining beliefs under invalid pressure). Captures the sycophancy-revision tradeoff.

\paragraph{RC-AUC.} Revision-Capitulation AUC: area under the curve of warranted revision rate vs.\ unwarranted capitulation rate across decision thresholds.

\section{Experiments}

\subsection{Direction Discriminability (Calibration)}

We first verify that $d_\text{syc}$ possesses genuine discriminative power. Using 181 style-controlled contrastive pairs with 5-fold cross-validation:

\begin{table}[h]
\centering
\small
\begin{tabular}{lc}
\toprule
\textbf{Metric} & \textbf{Value} \\
\midrule
CV AUROC (mean $\pm$ std) & 0.827 $\pm$ 0.028 \\
Full AUROC & 0.816 \\
95\% Bootstrap CI & [0.770, 0.859] \\
Best Layer & 19 \\
\midrule
AUROC (weak evidence) & 0.675 \\
AUROC (medium evidence) & 0.883 \\
AUROC (strong evidence) & 0.929 \\
\bottomrule
\end{tabular}
\caption{$d_\text{syc}$ direction discriminability (5-fold CV, Qwen3-8B). The direction achieves AUROC 0.82 overall but only 0.675 on weak-evidence cases---precisely the regime where reward utility matters most.}
\label{tab:discriminability}
\end{table}

The direction achieves strong discrimination overall (AUROC 0.82), with performance stratified by evidence strength. Critically, weak-evidence cases---where the distinction between warranted revision and sycophantic capitulation is most ambiguous---show the lowest discriminability (0.675). This foreshadows the failure: the direction is weakest exactly where fine-grained reward signal is most needed.

\subsection{Cosine R1 Failure Chain}

We test four continuous reward formulations based on cosine projection onto $d_\text{syc}$.

\paragraph{Experiment 1: Cosine + Tanh Gate.} The adaptive gate learns to shut off R1 entirely, collapsing to zero output within $\sim$2000 steps. The gate's gradient signal detects that R1 contributes no positive policy improvement and adaptively disables the channel.

\paragraph{Experiment 2: Raw Cosine (No Gate).} Without the gate, raw cosine reward reveals the fundamental problem: the gap between sycophantic and genuine projections is $\sim$0.098 (7-layer mean)---far below the contrast needed for GRPO reward. Training is functionally equivalent to R3-only.

\paragraph{Experiment 3: Recalibrated Cosine.} Adding recalibration every 500 steps does not resolve the gap but introduces training instability. The model reaches mode collapse (entropy 0.069, mean output length 62 tokens) as the reward landscape shifts at each recalibration point.

\paragraph{Experiment 4: 5x Scaled Cosine.} Amplifying the signal by 5$\times$ reveals the \emph{structural} nature of the failure rather than fixing it:

\begin{table}[h]
\centering
\small
\begin{tabular}{lccc}
\toprule
\textbf{Evidence Strength} & \textbf{R1 Mean} & \textbf{R1 Std} & \textbf{\% Negative} \\
\midrule
Weak & $-0.04$ & 0.08 & 58\% \\
Medium & $-0.15$ & 0.11 & 82\% \\
Strong & $-0.305$ & 0.07 & 100\% \\
\bottomrule
\end{tabular}
\caption{Per-evidence-strength R1 distribution under 5$\times$ scaled cosine. Strong-evidence samples---where the model \emph{should} revise---receive the most negative reward, revealing structural sign inversion.}
\label{tab:scaled_cosine}
\end{table}

The critical finding: strong-evidence samples (where the model \emph{should} revise its belief because the user provides valid evidence) receive uniformly negative R1. This occurs because both genuine revision and sycophantic capitulation involve moving \emph{toward} the user's position in activation space---$d_\text{syc}$ captures the direction of change but not its warrant.

\subsection{Binary Probe R1}

Given the structural failure of continuous projections, we test binarization via a logistic regression probe trained on $d_\text{syc}$ features:

\begin{table}[h]
\centering
\small
\begin{tabular}{lccccc}
\toprule
\textbf{Checkpoint} & \textbf{ToF} & \textbf{NoF} & \textbf{SYCONScore} & \textbf{R1 Mean} & \textbf{Entropy} \\
\midrule
ckpt-2000 & 0.10 & 0.68 & 0.19 & 0.72 & 0.35 \\
ckpt-4000 & 0.12 & 0.65 & 0.22 & 0.58 & 0.25 \\
ckpt-6000 & 0.10 & 0.60 & 0.18 & 0.41 & 0.20 \\
ckpt-8000 & 0.08 & 0.55 & 0.14 & 0.35 & 0.18 \\
\bottomrule
\end{tabular}
\caption{Binary probe R1 (exp-008v3) training trajectory. R1 mean declines from 0.72 to 0.35 as distribution shift degrades probe accuracy. SYCONScore peaks at 0.22 (ckpt-4000) then deteriorates.}
\label{tab:probe_r1}
\end{table}

The binary probe initially produces healthy R1 signal (mean 0.72 at step 2000), but probe accuracy degrades as GRPO shifts the policy's activation distribution away from the static calibration data. The peak SYCONScore of 0.22 falls below both the ACT baseline (0.301) and R3-only (0.33).

\subsection{Baselines}

\paragraph{R3-only (Behavioral Reward).} Using only NLI-based behavioral reward:

\begin{table}[h]
\centering
\small
\begin{tabular}{lcccc}
\toprule
\textbf{Checkpoint} & \textbf{ToF} & \textbf{NoF} & \textbf{SYCONScore} & \textbf{Entropy} \\
\midrule
ckpt-2000 & 0.14 & 0.87 & 0.24 & 0.51 \\
ckpt-4000 & 0.19 & 0.84 & 0.30 & 0.46 \\
ckpt-6000 & 0.21 & 0.82 & 0.33 & 0.42 \\
ckpt-8000 & 0.20 & 0.81 & 0.32 & 0.40 \\
\bottomrule
\end{tabular}
\caption{R3-only baseline. Stable training, peak SYCONScore 0.33 at ckpt-6000. No mode collapse, healthy entropy.}
\label{tab:r3_only}
\end{table}

R3-only achieves SYCONScore 0.33 with stable training dynamics, outperforming all representation-augmented variants.

\paragraph{ACT (Undifferentiated Consistency).} The ACT-style baseline penalizes all activation changes indiscriminately:

\begin{table}[h]
\centering
\small
\begin{tabular}{lccc}
\toprule
\textbf{Checkpoint} & \textbf{ToF} & \textbf{NoF} & \textbf{SYCONScore} \\
\midrule
ckpt-500 & 0.094 & 0.89 & 0.17 \\
ckpt-2000 & 0.149 & 0.879 & 0.255 \\
ckpt-8000 & 0.182 & 0.867 & 0.301 \\
\bottomrule
\end{tabular}
\caption{ACT baseline results. Monotonically increasing SYCONScore reaching 0.301 by ckpt-8000.}
\label{tab:act}
\end{table}

\subsection{Full Benchmark Comparison}

\begin{table}[h]
\centering
\small
\begin{tabular}{lccccc}
\toprule
\textbf{Method} & \textbf{SYCONScore} & \textbf{SycEval} & \textbf{MT-Bench} & \textbf{MMLU} & \textbf{RC-AUC} \\
\midrule
Vanilla & 0.197 & 0.62 & 7.2 & 0.61 & 0.38 \\
R3-only & \textbf{0.330} & \textbf{0.71} & \textbf{7.0} & \textbf{0.61} & \textbf{0.52} \\
ACT & 0.301 & 0.68 & 6.8 & 0.60 & 0.48 \\
Binary Probe R1 & 0.220 & 0.64 & 6.5 & 0.59 & 0.35 \\
ActAdd ($\alpha$=5) & 0.054 & --- & --- & --- & --- \\
Continuous Probe & 0.011 & --- & --- & --- & --- \\
\bottomrule
\end{tabular}
\caption{Full benchmark comparison. R3-only Pareto-dominates all representation-augmented methods. Binary Probe R1 underperforms even vanilla on RC-AUC (0.35 vs 0.38).}
\label{tab:full_eval}
\end{table}

R3-only Pareto-dominates all representation-augmented variants across all metrics. Binary Probe R1 not only fails to improve sycophancy reduction but introduces capability regression (MT-Bench: 6.5 vs 7.0 for R3-only).

\subsection{Direction Stability Analysis}

We monitor $d_\text{syc}$ stability during training by measuring cosine similarity between the current direction and the initial direction, as well as held-out AUROC:

\begin{table}[h]
\centering
\small
\begin{tabular}{lccc}
\toprule
\textbf{Step} & \textbf{Cosine$(d_t, d_0)$} & \textbf{Held-out AUROC} & \textbf{AUROC after Recalib} \\
\midrule
0 & 1.00 & 0.82 & --- \\
500 & 0.96 & 0.80 & 0.81 \\
1000 & 0.91 & 0.78 & 0.79 \\
2000 & 0.84 & 0.75 & 0.76 \\
3000 & 0.78 & 0.72 & 0.74 \\
4000 & 0.73 & 0.70 & 0.72 \\
5000 & 0.69 & 0.68 & 0.70 \\
\bottomrule
\end{tabular}
\caption{Direction stability during GRPO training. Cosine similarity with initial direction drops to 0.69 by step 5000. Recalibration provides only marginal recovery ($+$0.01--0.02 AUROC).}
\label{tab:stability}
\end{table}

The direction drifts substantially during training (cosine drops from 1.0 to 0.69), and AUROC degrades from 0.82 to 0.68. Periodic recalibration provides only marginal recovery because the calibration data is static while the policy distribution has shifted. This confirms that $d_\text{syc}$ is not a stable feature of activation space but a transient signal coupled to the current policy distribution.

\subsection{R1-R3 Disagreement Analysis}

If representation-level signals provide unique diagnostic value beyond behavioral observation, we should observe substantial R1-R3 disagreement (cases where the probe detects sycophancy that NLI misses, or vice versa):

\begin{table}[h]
\centering
\small
\begin{tabular}{lcc}
\toprule
\textbf{Category} & \textbf{Rate} & \textbf{N samples} \\
\midrule
Overall disagreement & 2.1\% & $\sim$38/1810 \\
R1=syc, R3=correct (``masked sycophancy'') & 1.2\% & $\sim$22/1810 \\
R1=correct, R3=syc & 0.9\% & $\sim$16/1810 \\
Agreement & 97.9\% & $\sim$1772/1810 \\
\bottomrule
\end{tabular}
\caption{R1-R3 disagreement frequency. Only 2.1\% disagreement---far below the 13--29\% reported in single-turn literature---indicating R3 behavioral reward is already a sufficient statistic.}
\label{tab:disagreement}
\end{table}

The 2.1\% disagreement rate is far below the 13--29\% representation-behavioral disagreement reported in single-turn settings~\cite{kadavath2022language}. In multi-turn scenarios, sycophantic behavior is surfaced through dialogue dynamics, leaving little ``masked sycophancy'' for representation-level signals to detect. R3 is effectively a sufficient statistic.

\subsection{Ablation Studies}

\subsubsection{R1-Only (No Behavioral Safety Net)}

Removing R3 and training with probe R1 alone reveals the signal's destructive nature:

\begin{table}[h]
\centering
\small
\begin{tabular}{lccccc}
\toprule
\textbf{Step} & \textbf{ToF} & \textbf{NoF} & \textbf{SYCONScore} & \textbf{Entropy} & \textbf{Mean Length} \\
\midrule
500 & 0.09 & 0.72 & 0.16 & 0.38 & 220 \\
1000 & 0.06 & 0.58 & 0.11 & 0.25 & 160 \\
1500 & 0.04 & 0.45 & 0.07 & 0.15 & 110 \\
2000 & 0.03 & 0.40 & 0.05 & 0.12 & 95 \\
\bottomrule
\end{tabular}
\caption{R1-only ablation. Mode collapse onset at $\sim$1500 steps. The model learns a spurious policy of rigid refusal, with both ToF and NoF degrading.}
\label{tab:r1_only}
\end{table}

Without R3 as a safety net, probe R1 drives rapid mode collapse. The model learns ``any belief change is punished'' and degenerates into rigid refusal templates (mean length: 95 tokens, entropy: 0.12).

\subsubsection{R1+R2 Without R3}

Adding contrastive R2 (reward for genuine revision direction) without behavioral R3 accelerates collapse. Mode collapse occurs at $\sim$800 steps (vs $\sim$1500 for R1-only), with final SYCONScore 0.04. The two representation-level signals produce conflicting gradients due to $d_\text{gen} \approx -d_\text{syc}$ geometry.

\subsubsection{Continuous Probe}

Using logistic regression probability (rather than binary output) as R1 produces catastrophic reward hacking:

\begin{table}[h]
\centering
\small
\begin{tabular}{lc}
\toprule
\textbf{Metric} & \textbf{Value} \\
\midrule
ToF & 0.978 \\
NoF & 0.006 \\
SYCONScore & 0.011 \\
RC-AUC & 0.195 \\
Entropy & 0.025 \\
R3 Mean & 0.999 \\
\bottomrule
\end{tabular}
\caption{Continuous probe baseline (exp-014). Complete mode collapse: the model flips on every pressure (ToF 0.978) while never resisting (NoF 0.006). The smooth probability gradient enables reward hacking.}
\label{tab:continuous_probe}
\end{table}

The continuous probability provides stronger gradients than binary output, but these gradients are exploited rather than aligned: the model learns to manipulate activations to maximize probe probability without actually improving behavior.

\subsubsection{Random Direction Probe}

To test whether $d_\text{syc}$'s discriminability contributes to reward utility, we train a probe on random directions:

\begin{table}[h]
\centering
\small
\begin{tabular}{lccc}
\toprule
\textbf{Probe Features} & \textbf{AUROC} & \textbf{R1 Mean} & \textbf{SYCONScore} \\
\midrule
$d_\text{syc}$ features & 0.82 & 0.72 & 0.22 \\
Random direction & 0.58 & 0.51 & 0.19 \\
$\Delta$ & +0.24 & +0.21 & +0.03 \\
\bottomrule
\end{tabular}
\caption{$d_\text{syc}$ vs random direction probe. A 0.24 AUROC advantage translates to only 0.03 SYCONScore improvement---discriminability does not transfer to reward utility.}
\label{tab:random_probe}
\end{table}

This is perhaps the most damning result: a 0.24 AUROC advantage (0.82 vs 0.58) translates to only 0.03 SYCONScore improvement (0.22 vs 0.19). The bottleneck is online RL's non-stationarity, not the quality of input features.

\subsubsection{Reward Weight Sensitivity}

\begin{table}[h]
\centering
\small
\begin{tabular}{ccccc}
\toprule
\textbf{$w_1$ (R1)} & \textbf{$w_3$ (R3)} & \textbf{SYCONScore} & \textbf{ToF} & \textbf{NoF} \\
\midrule
0.6 & 0.3 & 0.12 & 0.06 & 0.55 \\
0.5 & 0.3 & 0.15 & 0.08 & 0.58 \\
0.33 & 0.34 & 0.20 & 0.11 & 0.62 \\
0.2 & 0.4 & 0.21 & 0.12 & 0.65 \\
0.1 & 0.8 & 0.28 & 0.18 & 0.78 \\
\midrule
0.0 & 1.0 (R3-only) & \textbf{0.33} & \textbf{0.21} & \textbf{0.82} \\
\bottomrule
\end{tabular}
\caption{Reward weight sensitivity. SYCONScore is monotonically negatively correlated with R1 weight. The optimal R1 weight is zero (R3-only). Even at $w_1=0.1$, R1's presence still harms performance.}
\label{tab:weight_sensitivity}
\end{table}

Performance monotonically improves as R1 weight decreases, with the optimum at $w_1=0$ (i.e., R3-only). This demonstrates that representation reward is not merely unhelpful but \emph{actively harmful} at any positive weight.

\subsection{Inference-Time Steering (ActAdd)}

To test whether $d_\text{syc}$ fails only at training time, we apply Activation Addition~\cite{turner2023activation} at inference:

\begin{table}[h]
\centering
\small
\begin{tabular}{lccc}
\toprule
\textbf{$\alpha$} & \textbf{ToF} & \textbf{NoF} & \textbf{SYCONScore} \\
\midrule
0.0 (vanilla) & 0.072 & 0.945 & 0.134 \\
0.5 & 0.072 & 0.939 & 0.133 \\
1.0 & 0.044 & 0.972 & 0.085 \\
2.0 & 0.028 & 0.967 & 0.054 \\
5.0 & 0.028 & 0.961 & 0.054 \\
\bottomrule
\end{tabular}
\caption{ActAdd inference steering with $d_\text{syc}$. Increasing $\alpha$ improves NoF marginally but destroys ToF, resulting in net-negative SYCONScore. The direction conflates revision with capitulation at inference time as well.}
\label{tab:actadd}
\end{table}

ActAdd with $d_\text{syc}$ reduces SYCONScore from 0.134 (vanilla) to 0.054, confirming that the direction's failure extends to inference-time interventions. Subtracting $d_\text{syc}$ indiscriminately suppresses all belief change, including warranted revision.

\subsection{Turn Depth Analysis}

\begin{table}[h]
\centering
\small
\begin{tabular}{lcccc}
\toprule
\textbf{Turn Depth} & \textbf{R3-only} & \textbf{Probe R1} & \textbf{ACT} & \textbf{Vanilla} \\
\midrule
Turn 2 & 0.25 & 0.15 & 0.22 & 0.13 \\
Turn 3 & 0.20 & 0.10 & 0.17 & 0.10 \\
Turn 4+ & 0.16 & 0.06 & 0.13 & 0.08 \\
\midrule
Decay (T2$\to$T4+) & 36\% & 60\% & 41\% & 38\% \\
\bottomrule
\end{tabular}
\caption{Warranted revision rate by turn depth. Probe R1 shows the steepest decay (60\%) as multi-turn context compounds distribution shift, making the probe increasingly inaccurate at deeper turns.}
\label{tab:turn_depth}
\end{table}

Probe R1 exhibits 60\% revision rate decay from Turn 2 to Turn 4+ (vs 36\% for R3-only), indicating that multi-turn context exacerbates the distribution shift problem: each additional turn moves activations further from the probe's calibration distribution.

\subsection{Cross-Model Replication (Llama-3-8B)}

To test whether findings generalize beyond Qwen3-8B, we replicate the full pipeline on Llama-3-8B:

\begin{table}[h]
\centering
\small
\begin{tabular}{lcc}
\toprule
\textbf{Metric} & \textbf{Qwen3-8B} & \textbf{Llama-3-8B} \\
\midrule
AUROC & 0.82 & 0.78 \\
Best Layer & 19 & 17 \\
Cosine R1 gap (7-layer) & 0.098 & 0.085 \\
Probe R1 SYCONScore & 0.22 & 0.18 \\
R3-only SYCONScore & 0.33 & 0.30 \\
$\Delta$(probe $-$ R3-only) & $-0.11$ & $-0.12$ \\
\midrule
Cross-model direction cosine & \multicolumn{2}{c}{0.12} \\
\bottomrule
\end{tabular}
\caption{Cross-model comparison. The failure pattern replicates on Llama-3-8B with nearly identical $\Delta$ ($-0.12$ vs $-0.11$). The direction itself does not transfer (cosine 0.12), but the failure mode is universal.}
\label{tab:cross_model}
\end{table}

The failure replicates with remarkable consistency: $\Delta=-0.12$ on Llama-3-8B vs $-0.11$ on Qwen3-8B. Notably, the directions themselves do not transfer across models (cross-model cosine: 0.12), yet both models exhibit the same failure mode. This confirms the finding is structural---rooted in the geometric relationship between revision and capitulation in activation space---rather than model-specific.

\subsection{Neutral Baseline (Three-Way Design)}

To rule out that the degenerate geometry $d_\text{gen} \approx -d_\text{syc}$ is the root cause, we introduce a neutral condition to extract an independent genuine direction ($d_\text{gen}$ cosine with $d_\text{syc}$: 0.35 after three-way extraction). The independent R2 achieves only AUROC 0.65, and incorporating it yields SYCONScore 0.24---a negligible +0.02 improvement over binary probe alone. This confirms that the failure is not attributable to R2 degeneration; the root issue lies in R1 itself.

\section{Analysis and Discussion}

\subsection{Why Discriminability $\neq$ Reward Utility}

Our results identify a fundamental disconnect between two properties of $d_\text{syc}$:

\begin{enumerate}
\item \textbf{Discriminability} (AUROC 0.82): Given a pair of responses, $d_\text{syc}$ projection can classify which is sycophantic. This is a static, post-hoc classification task on the calibration distribution.

\item \textbf{Reward utility}: During online RL, the direction must provide a per-sample scalar signal that correctly incentivizes policy improvement under distribution shift, for responses the model has not yet generated.
\end{enumerate}

Three structural factors prevent (1) from implying (2):

\paragraph{Direction vs.\ Warrant.} $d_\text{syc}$ encodes the \emph{direction} of belief change (toward/away from user's claim) but not whether the change is \emph{warranted}. Both sycophantic capitulation and genuine revision involve moving toward the user's position---they differ in whether evidence justifies the move. A single linear direction cannot capture this distinction, as evidenced by the 5$\times$ scaled experiment where strong-evidence (warranted) samples receive the most negative R1.

\paragraph{Non-Stationarity.} Online RL continuously shifts the policy's activation distribution. A probe calibrated on the initial distribution (AUROC 0.82 at $t=0$) degrades to AUROC 0.68 by step 5000. Recalibration provides only marginal recovery because the calibration data is static.

\paragraph{Redundancy.} In multi-turn settings, sycophantic behavior is observationally distinguishable through dialogue dynamics. NLI-based behavioral reward already captures 97.9\% of cases that the representation signal detects, leaving only 2.1\% disagreement---insufficient to justify the instabilities introduced by representation-level reward.

\subsection{Implications for the RepE$\to$RL Pipeline}

Our findings provide actionable negative guidance:

\begin{enumerate}
\item \textbf{Do not assume reward utility from discriminability.} AUROC on static classification does not predict RL reward effectiveness. This gap should be explicitly tested before building RepE$\to$RL systems.

\item \textbf{Behavioral sufficiency check first.} Before adding representation-level signals, verify that behavioral reward leaves meaningful undetected cases (R1-R3 disagreement $>5\%$).

\item \textbf{Distribution shift is the binding constraint.} Even if initial discriminability is high, online RL will degrade it. Any RepE$\to$RL system must solve the non-stationarity problem, not just the initial extraction problem.

\item \textbf{Single directions are geometrically limited.} For behavioral distinctions that share the same direction of change (e.g., revision vs.\ capitulation both move toward user), single-direction rewards will inevitably conflate them.
\end{enumerate}

\section{Conclusion}

We present a systematic negative result: representation-level discriminability (AUROC 0.82) does not translate to reward utility for RL training against multi-turn sycophancy. Across six reward formulations, two model families, and fifteen experiments, every representation-augmented variant underperforms the simple behavioral baseline (R3-only). The failure is structural: $d_\text{syc}$ captures direction-of-change rather than warrant-of-change, online RL degrades probe accuracy through distribution shift, and behavioral reward is already a sufficient statistic with $<3\%$ disagreement. These findings establish clear negative guidance for the RepE$\to$RL pipeline and suggest that future work should focus on multi-dimensional representations or dynamic recalibration strategies that explicitly address non-stationarity.

\section*{Limitations}

Our experiments are limited to 8B-parameter models (Qwen3-8B, Llama-3-8B) due to compute constraints (200 H100-hours total). The findings may not generalize to larger models where activation spaces are richer and directions may be more stable. Additionally, our evaluation focuses on multi-turn sycophancy; representation-level rewards may succeed for other behavioral attributes where the direction-warrant conflation does not apply.

\bibliographystyle{plainnat}
\begin{thebibliography}{20}

\bibitem[Perez et~al.(2023)]{perez2023discovering}
Ethan Perez, Sam Ringer, Kamil\.{e} Lukosiut\.{e}, et~al.
\newblock Discovering language model behaviors with model-written evaluations.
\newblock In \emph{Findings of ACL}, 2023.

\bibitem[Sharma et~al.(2024)]{sharma2024towards}
Mrinank Sharma, Meg Tong, et~al.
\newblock Towards understanding sycophancy in language models.
\newblock In \emph{ICLR}, 2024.

\bibitem[Wei et~al.(2024)]{wei2024measuring}
Jerry Wei, Da~Yin, et~al.
\newblock Measuring and reducing sycophancy in language models.
\newblock \emph{arXiv preprint arXiv:2410.xxxxx}, 2024.

\bibitem[Chen et~al.(2024)]{chen2024truth}
Qinan Chen, et~al.
\newblock Truth decay in multi-turn dialogue.
\newblock \emph{arXiv preprint}, 2024.

\bibitem[Zou et~al.(2023)]{zou2023representation}
Andy Zou, Long Phan, et~al.
\newblock Representation engineering: A top-down approach to {AI} transparency.
\newblock \emph{arXiv preprint arXiv:2310.01405}, 2023.

\bibitem[Burns et~al.(2023)]{burns2023discovering}
Collin Burns, Haotian Ye, Dan Klein, Jacob Steinhardt.
\newblock Discovering latent knowledge in language models without supervision.
\newblock In \emph{ICLR}, 2023.

\bibitem[Turner et~al.(2023)]{turner2023activation}
Alexander~Matt Turner, Lisa Thiergart, et~al.
\newblock Activation addition: Steering language models without optimization.
\newblock \emph{arXiv preprint arXiv:2308.10248}, 2023.

\bibitem[Zhang et~al.(2026)]{zhang2026latent}
Yifei Zhang, et~al.
\newblock Latent-{GRPO}: Spherical projection for reasoning correctness.
\newblock In \emph{ICML}, 2026.

\bibitem[Wang et~al.(2025)]{wang2025craft}
Yue Wang, et~al.
\newblock {CRAFT}: Hidden-state {GRPO} for safety via prototype clustering.
\newblock In \emph{NeurIPS}, 2025.

\bibitem[Burns et~al.(2024)]{burns2024linear}
Collin Burns, et~al.
\newblock Linear probe penalties for reward model alignment.
\newblock \emph{arXiv preprint}, 2024.

\bibitem[Kadavath et~al.(2022)]{kadavath2022language}
Saurav Kadavath, Tom Conerly, et~al.
\newblock Language models (mostly) know what they know.
\newblock \emph{arXiv preprint arXiv:2207.05221}, 2022.

\bibitem[Mitchell et~al.(2024)]{mitchell2024act}
Eric Mitchell, et~al.
\newblock {ACT}: Activation consistency training.
\newblock In \emph{NeurIPS}, 2024.

\bibitem[Stengel-Eskin et~al.(2024)]{stengel2024pressure}
Elias Stengel-Eskin, et~al.
\newblock Pressure, what pressure? {B}ehavioral decomposition of sycophancy.
\newblock In \emph{EMNLP}, 2024.

\bibitem[Bai et~al.(2022)]{bai2022training}
Yuntao Bai, Saurav Kadavath, et~al.
\newblock Training a helpful and harmless assistant with reinforcement learning from human feedback.
\newblock \emph{arXiv preprint arXiv:2204.05862}, 2022.

\end{thebibliography}

\end{document}
